\subsection{BASES OF A FILTER}

If \(\mathfrak{G}\) is a subbase of a filter \(\mathfrak{F}\) on X (no. 2), then \(\mathfrak{F}\) is not in general the set of subsets of X which contain a set of \(\mathfrak{G}\) ; for \(\mathfrak{G}\) to have this property it is necessary and sufficient that every finite intersection of sets of \(\mathfrak{G}\) should contain a set of \(\mathfrak{G}\) . Hence the following proposition:

PROPOSITION 2. Let \(\mathfrak{B}\) be a set of subsets of a set X. Then the set of subsets of X which contain a set of \(\mathfrak{B}\) is a filter if and only if \(\mathfrak{B}\) has the following two properties:

(B \(_{I}\) ) The intersection of two sets of B contains a set of B.

(B \(_{II}\) ) B is not empty, and the empty subset of X is not in B.

DEFINITION 3. A set B of subsets of a set X which satisfies axioms (B \(_{I}\) ) and (B \(_{II}\) ) is said to be a base of the filter it generates. Two filter bases are said to be equivalent if they generate the same filter.

If \(\mathfrak{G}\) is a subbase of a filter \(\mathfrak{F}\) , then the set \(\mathfrak{G}'\) of finite intersections of sets of \(\mathfrak{G}\) is a base of \(\mathfrak{F}\) (no. 2).

PROPOSITION 3. A subset \(\mathfrak{B}\) of a filter \(\mathfrak{F}\) on \(\mathbf{X}\) is a base of \(\mathfrak{F}\) if and only if every set of \(\mathfrak{F}\) contains a set of \(\mathfrak{B}\) .

If \(\mathfrak{B}\) is a base of \(\mathfrak{F}\) , then clearly every set of \(\mathfrak{F}\) contains a set of \(\mathfrak{B}\) ; conversely, if every set of \(\mathfrak{F}\) contains a set of \(\mathfrak{B}\) , then the set of subsets of \(X\) containing a set of \(\mathfrak{B}\) coincides with \(\mathfrak{F}\) by reason of \((\mathbf{F}_{\mathbf{I}})\) .

For example, the Fréchet filter (no. 1) is the section filter of the ordered set N, considered as directed by the relation \(\leqslant\). Let \(\mathfrak{F}\) be a filter on a set Z. Since \(\mathfrak{F}\) is directed with respect to the relation \(\supset\) {[}by reason of axiom \((\mathrm{F}_{\mathrm{II}})\){]} we can define a section filter on \(\mathfrak{F}\); here a section of \(\mathfrak{F}\) relative to a set A\(\in\mathfrak{F}\) is the set S(A) of all M\(\in\mathfrak{F}\) such that M\(\subset\)A. This filter is called the section filter of the filter \(\mathfrak{F}\).

PROPOSITION 4. On a set X, a filter \(\mathfrak{F}'\) with base \(\mathfrak{B}'\) is finer than a filter \(\mathfrak{F}\) with base \(\mathfrak{B}\) if and only if every set of \(\mathfrak{B}\) contains a set of \(\mathfrak{B}'\) .

This is an immediate consequence of Definitions 2 and 3.

COROLLARY. Two filter bases B, B' on a set X are equivalent if and only if every set of B contains a set of B' and every set of B' contains a set of B.

Examples of filter bases. 1) Let X be a topological space. Proposition 3 shows that the bases of the neighbourhood filter of a point \(x \in X\) are precisely the fundamental systems of neighbourhoods of \(x\) (§ 1, no. 3, Definition 5).

\begin{enumerate}
\def\labelenumi{\arabic{enumi})}
\setcounter{enumi}{1}
\tightlist
\item
  Let X be a non-empty directed set with respect to a relation \((\sigma)\) (Set Theory, Chapter III, § 1, no. 10). For each \(a \in X\) , the set \(S(a)\) of all \(x \in X\) such that \(a(\sigma)x\) will be called the section of X relative to the element a. Then the set S of sections of X is a filter base, for it clearly satisfies \((\mathbf{B}_{\mathbf{II}})\) , and if a, b are any two elements of X, then there is by hypothesis an element \(c \in X\) such that \(a(\sigma)c\) and \(b(\sigma)c\) , and therefore
\end{enumerate}

\[
\mathbf {S} (c) \subset \mathbf {S} (a) \cap \mathbf {S} (b),
\]

so that \((\mathbf{B}_{\mathbf{I}})\) is satisfied. The filter generated by S is called the section filter of the directed set X.

